\scanpage{145}
Si $E$ est un Banach, et $u_\alpha\to u$ \unsure{dans|ds} $L_0(E)$, alors $\sigma(u_\alpha)\to\sigma(u)$ dans le sens suivant\unsure{.|:}

De plus, convergence des projecteurs spectraux \unsure{--\,--\,--\,--\,--\, .} Dans le cas des op. hermitiens dans l'espa\unsure{ce} de Hilbert, on a $=$ \unsure{là|la} des conditions nécessaires et suffisantes de convergence.

Convergence de \unsure{l'hermitien|d'hermitiens|l'hermiticité} \inserted{$u_\alpha\to u$} dans $L^p(E)$ : il est nécessaire et suffisant que \textcircled{\scriptsize 1°} $\|u_\alpha\|_p\to\|u\|_p$, \textcircled{\scriptsize 2°} $u_\alpha\to u$ spectralement, i.e, dans $L(E)$

(il suffit \unsure{$=$|même} que l'on ait la convergence spectrale individuelle en chaque $\lambda\in\sigma(u)$, $\lambda\neq0$)

\uline{Corollaire 1} $\unsure{U|u}\to|A|$, de $L^p(E,F)$ dans $L^p(E)$ est continue. Plus généralement, $u\to|u|^r$ de $L^p(E,F)$ dans $L^{p/r}(E)$ est continue.

Car composée de $u\to u^*u$ de $L^p(E,F)$ dans $(L^{p/2}(E))^+$, avec $v\to v^{\frac r2}$ de $(L^{p/2}(E))^+$ dans $L^{\frac pr}(E)$ ; or \deleted{\unsure{la|cette}} \unsure{dernière|dernier} est continue, car elle est continue de $L_0(E)^+$ dans $L_0(E)$ (ce dernier \unsure{point|fait} se vérifie \uline{spectralement} ; c'est sans doute vrai aussi de $L(E)^+$ dans $L(E)$ ? Inégalités précises ?)

\uline{Corollaire 2} Si $0<p\leq\unsure{2|1}$, pour qu'une suite \inserted{d'hermitiens} $u_\alpha\to u$ converge dans $L^p(E\deleted{\unsure{,F|,E}})$, il faut et il suffit que $\|u_\alpha\|_p\to\|u\|_p$, et qu'on ait la convergence faible\unsure{.}
